\paragraph{Обозначения.} Через $\varepsilon>0$ обозначается требуемая точность решения задачи по функции: метод должен выдать точку $\hat x$, для которой $f(\hat x)-f^\ast\le\varepsilon$. Знак $\norm{\cdot}$ всюду означает евклидову норму вектора и спектральную норму матрицы, $[x,y]$ --- отрезок с концами $x,y\in\R^d$, $B(x,r)$ --- замкнутый шар радиуса $r$ с центром $x$. Сумму первых $k$ членов последовательности мы обозначаем соответствующей заглавной буквой: например, $A_k=a_1+\dots+a_k$ и $A_0=0$.
